\section[Literature Review]{Literature Review and State of the Art}
\label{sec:pd_literature}

The Ge-on-Si waveguide photodetector field has evolved over approximately two
decades through successive generations of electrode and junction engineering,
each interpretable through the metric framework of \cref{sec:pd_metrics}.
What follows is a causal narrative: each development is best understood as an
attempt to relax one bottleneck in \cref{eq:bw_total,eq:idark_decomposed,eq:qe_full}
while leaving another constraint behind.

\subsection[Milestone Waveguide-Integrated PIN Photodetectors]{Milestone
Waveguide-Integrated Ge-on-Si PIN Photodetectors}
\label{subsec:pd_history}

The landmark demonstration of Vivien~\textit{et al.}~\cite{Vivien2009} established
a silicon-on-insulator (SOI) waveguide-coupled Ge p-i-n detector with
$\mathcal{R}\approx\SI{1.0}{\ampere\per\watt}$ and
$f_{3\,\mathrm{dB}}=\SI{42}{\giga\hertz}$ at the O-band, confirming that
direct-gap absorption at the $\Gamma$-point translates into competitive
responsivity on the silicon platform. The \SI{42}{\giga\hertz} result was
RC-limited: the electrode geometry left $R_s$ at a value that caused
$f_{\mathrm{RC}}$ to dominate over $f_{\mathrm{tt}}$ in \cref{eq:bw_total}.
Ang~\textit{et al.}~\cite{Ang2010} subsequently demonstrated evanescently coupled
Ge-on-SOI fabricated at low thermal budget to remain CMOS-compatible, establishing
process integration as a design constraint alongside responsivity and bandwidth.

The next phase of the literature attacked $R_sC_j$ more directly through contact
and pad engineering. Chen~\textit{et al.}~\cite{Chen2015} demonstrated a
silicon-contacted Ge p-i-n with $\mathcal{R}=\SI{0.74}{\ampere\per\watt}$ and
$f_{3\,\mathrm{dB}}=\SI{35}{\giga\hertz}$, and the same group later pushed to
$\SI{67}{\giga\hertz}$ at only \SI{1}{\volt} reverse bias through electrode
repositioning~\cite{Chen2016Si}, thereby reducing $\ell_a$ in \cref{eq:rs_anode}
until $f_{\mathrm{RC}}$ no longer dominated \cref{eq:bw_total}.
Novack~\textit{et al.}~\cite{Novack2013} reached $\SI{60}{\giga\hertz}$ by
adding wire-bond inductance as a gain-peaking element. That approach is effective
but system-dependent: the added inductance partially resonates with $C_j+C_{\mathrm{par}}$,
so the bandwidth extension depends on the surrounding RF environment and package
parasitics~\cite{Yan2024}.

The lateral junction architecture introduced by Virot~\textit{et al.}~\cite{Virot2017}
decouples the optical mode depth from the depletion region width, providing an
additional degree of freedom for $C_j$ minimisation in \cref{eq:cj} that is
geometrically independent of $\mathcal{A}$ in \cref{eq:abs_fraction}. Parasitic
engineering of the vertical architecture was pushed to its prior limit by
Shi~\textit{et al.}~\cite{Shi2021}, who demonstrated
$f_{3\,\mathrm{dB}}=\SI{80}{\giga\hertz}$ through systematic minimisation of
parasitic pad capacitance. The bandwidth ceiling for the lateral architecture
was located by Lischke~\textit{et al.}~\cite{Lischke2021}, who reported
$f_{3\,\mathrm{dB}}=\SI{265}{\giga\hertz}$ with a Ge-fin geometry scaled to
$W_{\mathrm{Ge}}\approx\SI{100}{\nano\metre}$.

\paragraph{Resonant-cavity-enhanced and uni-travelling-carrier alternatives.}
Two architectural alternatives provide additional context. The
resonant-cavity-enhanced (RCE) vertical photodiode places the absorber within
a Fabry-P\'erot resonator; at the resonance condition, the single-pass absorbed
fraction $\mathcal{A}$ of \cref{eq:abs_fraction} is replaced by:
\begin{equation}
    \mathcal{A}_{\mathrm{RCE}}
    = \frac{(1-R_1)(1+R_2\,e^{-\alpha L})^2(1-e^{-2\alpha L})}
           {1-2\sqrt{R_1R_2}\,e^{-\alpha L}\cos(4\pi nL/\lambda)
            +R_1R_2\,e^{-2\alpha L}},
    \label{eq:rce_absorption}
\end{equation}
where $R_1$ and $R_2$ are the top and bottom mirror reflectances~\cite{Yan2024}.
Near-unity quantum efficiency is achievable in absorbers as thin as
\SI{200}{\nano\metre}, but the resonance bandwidth narrows to a few nanometres,
which is incompatible with the temperature-insensitive broadband operation required
of data-centre receivers. The uni-travelling-carrier (UTC) structure replaces
both drift species with electrons only, raising the transit-time limit to:
\begin{equation}
    f_{\mathrm{tt,UTC}}=\frac{0.45\,v_{\mathrm{d,sat},e}}{W_{\mathrm{dep}}},
    \label{eq:bw_tt_utc}
\end{equation}
approximately \SI{50}{\percent} higher than the hole-limited value of
\cref{eq:bw_tt} for the same depletion width~\cite{Yan2024}.

\subsection[Bandwidth Enhancement Strategies]{Bandwidth Enhancement Strategies}
\label{subsec:pd_bw_strategies}

The literature reveals four categories of bandwidth enhancement strategy:

\emph{Geometry scaling} reduces both $C_j$ (via smaller $A_j$) and $f_{\mathrm{tt}}$
(via smaller $W_i$), but doing so simultaneously sacrifices responsivity through
\cref{eq:qe_full}. Diminishing returns set in once the device enters the
absorption-starved regime below $L\approx\SI{4}{\micro\metre}$.

\emph{Doping profile optimisation} targets the dark-current-to-responsivity
balance: highly doped contacts reduce $R_s$ but increase free-carrier absorption,
while the junction polarity affects the heterojunction band alignment and
therefore $J_{\mathrm{diff}}$ through \cref{eq:jdark_diff}.

\emph{Contact engineering} attempts to reduce the dominant anode series resistance
through silicidation, via placement, and metal routing. For a vertical architecture
where $R_s$ is dominated by the lateral Si slab resistance of \cref{eq:rs_anode},
silicidation reduces the metal-semiconductor contact contribution but not the bulk
Si slab contribution, leaving $R_s$ improvement bounded by geometry.

\emph{Electrode topology} is the most structurally impactful lever: changing the
physical layout of the metal anode changes $\ell_a$ in \cref{eq:rs_anode} and
therefore $R_s$ directly. The U-shaped-anode strategy targets this lever by
reducing $\ell_a$ geometrically until $f_{\mathrm{RC}}>f_{\mathrm{tt}}$, providing
a structural solution rather than a resonant compensation.

\subsection[Advanced Electrode Designs Including U-Shaped Configurations]{Advanced
Electrode Designs Including U-Shaped Configurations}
\label{subsec:pd_u_electrode_review}

The benchmark for the present work is the \SI{103}{\giga\hertz} vertical Ge-on-Si
photodiode of Shi~\textit{et al.}~\cite{Shi2024}, which surpasses the
\SI{100}{\giga\hertz} threshold in the vertical architecture through a U-shaped
anode electrode that reduces $R_s$, and therefore $\tau_{\mathrm{RC}}=R_sC_j$ of
\cref{eq:tau_rc}, by approximately \SI{36}{\percent} relative to a
straight-electrode design. The U-shaped topology wraps the metal anode around
three sides of the germanium mesa, so that the maximum lateral current path from
any point on the mesa base to a metal contact is $W_m/4+d_{\mathrm{gap}}$ rather
than $W_m/2+d_{\mathrm{gap}}$. The U-shaped design simultaneously delivers
$\mathcal{R}=\SI{0.95}{\ampere\per\watt}$ and $I_d=\SI{1.3}{\nano\ampere}$ at
\SI{-1}{\volt}, making it a balanced rather than a specialised solution.

The U-shape is therefore selected not because it changes only a single resistor
value, but because it opens a favourable region of the coupled $(R_s,C_j,L_p)$
design space: the reduced $R_s$ also enables a shallower and more useful
inductive peaking condition, and the experimentally observed bandwidth increases
from \SI{83}{\giga\hertz} (parallel electrode) to \SI{103}{\giga\hertz}
(single-U)~\cite{Yan2024}.

\subsection[Performance Benchmarking]{Performance Benchmarking Table}
\label{subsec:pd_benchmarks}

\Cref{tab:pd_lit_benchmarks} consolidates the landmark experimental results most
relevant to the present design argument. The table is intentionally restricted to
the vertical trajectory adopted in this chapter and the lateral trajectory that
defines the broader bandwidth ceiling.

\begin{table}[H]
\centering
\caption{Selected benchmark waveguide-integrated Ge-on-Si photodetectors used to
position the present design. ``V'' and ``L'' denote vertical and lateral junction
geometries, respectively; ``NR'' denotes values not reported by the authors.}
\label{tab:pd_lit_benchmarks}
\footnotesize
\begin{tabular*}{\linewidth}{@{\extracolsep{\fill}}C{0.8cm}C{0.8cm}L{1.8cm}L{5.1cm}C{1.0cm}C{1.1cm}C{1.3cm}@{}}
\toprule
Year & Type & Reference & Key Design Lever &
$\mathcal{R}$ [A/W] &
$I_d$ [nA] &
$f_{3\,\mathrm{dB}}$ [GHz] \\
\midrule
2008 & V & \cite{Ang2010}
     & Evanescent coupling, CMOS-compatible process & 0.92 & 570 & 11.3 \\
2009 & V & \cite{Vivien2009}
     & Butt coupling, straight contact & $\sim$1.0 & 18 & 42 \\
2021 & V & \cite{Shi2021}
     & Parasitic engineering with inductive peaking & 0.89 & 6.4 & 80 \\
\textbf{2024} & \textbf{V} & \textbf{\cite{Shi2024}}
     & \textbf{U-shaped anode with inductive peaking} & \textbf{0.95} & \textbf{1.3} & \textbf{103} \\
2015 & L & \cite{Lischke2015}
     & Reduced Ge width with lateral offset junction & 1.00 & 100 & $>70$ \\
2017 & L & \cite{Virot2017}
     & Si/Ge/Si heterojunction with sub-micrometre Ge width & NR & 5 & $>50$ \\
2021 & L & \cite{Lischke2021}
     & Ge-fin junction with \SI{100}{\nano\metre} minimum width & 0.30 & 100 to 200 & \textbf{265} \\
2024 & L & \cite{Cao2024}
     & Side-coupled parallel Ge geometry & 0.80 & 24 & $>40$ \\
\bottomrule
\end{tabular*}
\end{table}

\begin{figure}[H]
    \centering
    \missingfigure[figwidth=0.88\linewidth]{
        Two-panel literature-trend figure. Left panel: vertical Ge-on-Si
        waveguide photodetectors plotted as bandwidth versus year, with markers
        distinguished by the main design lever. Right panel: lateral Ge-on-Si
        waveguide photodetectors on the same bandwidth-versus-year basis.
    }
    \caption{State-of-the-art trend summary for waveguide-integrated Ge-on-Si
    photodetectors. The chronology is most useful when read as a sequence of
    interventions into $R_s$, $C_j$, and the transport path rather than as a
    collection of isolated record values.}
\label{fig:pd_literature_trends}
\end{figure}

\subsection[Identified Gaps and Motivation]{Identified Gaps and Motivation for This Work}
\label{subsec:pd_gaps}

The identified gaps that motivate the present work are: (i) no prior published
simulation study has re-optimised the U-shaped vertical architecture for
\SI{1310}{\nano\metre} O-band operation, where the optical absorption and mode
coupling differ from the \SI{1550}{\nano\metre} benchmark; (ii) the double
U-shaped electrode topology, which wraps metal on both long sidewalls, remains
an open design concept rather than a validated implementation in the present
repository; and (iii) a reproducible solver-to-thesis workflow linking FDTD,
CHARGE, MATLAB, INTERCONNECT, and layout generation is rarely documented at
device-project level. These gaps frame the design and analysis scope documented
in \cref{sec:pd_design,sec:pd_simulation}.

\paragraph{IEEE~802.3 400GBASE-DR4 as the receiver-level target.}
IEEE~802.3 Clause~140 defines 400GBASE-DR4 as a \SI{400}{\giga\bit\per\second}
single-mode-fibre interface in the O-band, implemented through four parallel
optical lanes at \SI{1310}{\nano\metre} over reaches up to
\SI{500}{\metre}~\cite{IEEE80232022}. Each lane carries \SI{106.25}{\giga\bit\per\second}
raw line rate with PAM-4 signalling. The key receiver-side quantities are
summarised in \cref{tab:ieee_specs}.

\begin{table}[H]
\centering
\caption{IEEE~802.3 400GBASE-DR4 receiver-side constraints most directly relevant
to the photodetector, extracted from Clause~140~\cite{IEEE80232022} and mapped
onto the device metrics of \cref{sec:pd_metrics}.}
\label{tab:ieee_specs}
\footnotesize
\begin{tabular*}{\linewidth}{@{\extracolsep{\fill}}L{4.2cm}C{1.9cm}C{1.7cm}C{1.0cm}L{3.3cm}@{}}
\toprule
Parameter & Symbol & Value & Unit & Principal relation \\
\midrule
Wavelength range & $\lambda$ & 1260--1360 & nm & \cref{eq:abs_fraction} \\
Symbol rate & $f_s$ & 53.125 & Gbaud & \cref{eq:bw_total} \\
Modulation format & --- & PAM-4 & --- & --- \\
Minimum receiver sensitivity (OMA) & $\mathrm{OMA_{out,min}}$ & $-6.0$ & dBm & --- \\
Maximum input power & $P_{\mathrm{in,max}}$ & $+3.4$ & dBm & \cref{eq:p_screening} \\
Minimum receiver $-3\,\mathrm{dB}$ bandwidth & $f_{3\,\mathrm{dB,min}}$ & $\geq0.75f_s$ & GHz & \cref{eq:bw_total} \\
Reflectance tolerance & $\Gamma_{\mathrm{max}}$ & $-26$ & dB & \cref{eq:fdtd_abs} \\
\bottomrule
\end{tabular*}
\end{table}

The minimum receiver bandwidth implied by the conventional $0.75f_s$ criterion
is only \SI{39.8}{\giga\hertz}, so the benchmark vertical topology already
establishes comfortable speed relevance for DR4. The selected U-shaped vertical
Ge-on-Si PIN of Shi~\textit{et al.}~\cite{Shi2024} combines \SI{103}{\giga\hertz}-class
response, $\mathcal{R}=\SI{0.95}{\ampere\per\watt}$, and
$I_d=\SI{1.3}{\nano\ampere}$ in a process-compatible O-band platform.
Because DR4 requires a detector that is simultaneously fast enough, efficient at
low OMA, and quiet enough not to erode receiver margin, this balanced vertical
U-shaped architecture is the most defensible starting point for the present design.

\paragraph{Microhole enhancements and DBR implementations.}
Two structural variants highlight alternative optimization paths for bandwidth
and responsivity. To decouple the absorption length from the transit time,
Hamadou~\textit{et al.}~\cite{Hamadou2023} investigated surface-illuminated
silicon PIN photodiodes structured with photon-trapping microholes.
Their 3D Lumerical FDTD and CHARGE simulations demonstrated that etching a
$2\times2$ array of \SI{0.7}{\micro\metre} square microholes into the intrinsic
region enhances optical absorption through photon trapping while simultaneously
reducing the junction capacitance by a geometric filling factor $\rho_f$.
By carefully minimizing parasitic bonding capacitance ($C_b=\SI{70}{\femto\farad}$)
and series resistance ($R_s=\SI{3.6}{\ohm}$), their
\SI{20}{\micro\metre}$\times$\SI{20}{\micro\metre} device
achieved an open PAM-4 eye at \SI{50}{\giga\baud}.
Separately, the seminal work of Li~\textit{et al.}~\cite{Li2012} demonstrated
a \SI{112}{\giga\bit\per\second} CMOS-compatible waveguide Germanium photodetector
by utilizing a distributed Bragg reflector (DBR). The DBR allowed the optical
absorption length to be doubled without extending the physical transport distance,
while their interdigitated electrode geometry highlighted the delicate balance
between transit-time reduction and the onset of high dark current from strong
electric fields.

Broader reviews of the design space are provided by Wang~\textit{et al.}~\cite{Wang2024}
and Yan~\textit{et al.}~\cite{Yan2024}; the multiphysics modelling context is
treated in depth in~\cite{Alasio2024}; computational methods are discussed in
the pedagogical framework of Wartak~\cite{Wartak2013}.

\subsection[Carrier Transport Fundamentals in Reverse-Biased PIN Junctions]
{Carrier Transport Fundamentals in Reverse-Biased PIN Junctions}
\label{subsec:pd_transport_fundamentals}

A rigorous understanding of photodetector bandwidth and noise requires
first establishing the microscopic carrier dynamics within the depleted
intrinsic region of a reverse-biased p-i-n junction. What follows
develops the drift-diffusion framework from first principles, interprets
the resulting transit-time and RC-limited bandwidth expressions
physically, and draws direct consequences for the vertical Ge-on-Si
architecture operating at \SI{1310}{\nano\metre}.

\subsubsection{Carrier Generation and the Photocurrent Mechanism}

When a photon of energy $h\nu > E_g$ is absorbed in the intrinsic
germanium layer, a bound electron-hole pair is promoted across the
direct bandgap at the $\Gamma$-point of the Brillouin zone. For
germanium at the O-band centre wavelength of \SI{1310}{\nano\metre},
the photon energy is $h\nu \approx \SI{0.946}{\electronvolt}$, which
exceeds the direct-gap energy $E_{\Gamma} \approx \SI{0.80}{\electronvolt}$
at room temperature. This direct-gap absorption mechanism is responsible
for the significantly higher absorption coefficient of germanium relative
to silicon at this wavelength; whereas silicon exhibits
$\alpha_{\mathrm{Si}} \lesssim \SI{1e3}{\per\metre}$ at \SI{1310}{\nano\metre},
germanium achieves $\alpha_{\mathrm{Ge}} \approx \SI{4e4}{\per\metre}$
under the tensile strain conditions typical of CMOS-compatible epitaxial
growth on silicon~\cite{Vivien2009}. The implication for device design
is immediate: a Ge absorber of length $L \approx \SI{10}{\micro\metre}$
in a waveguide geometry can approach the single-pass absorption fraction
$\mathcal{A} \rightarrow 1 - e^{-\alpha_{\mathrm{Ge}} L} \approx 0.33$
even without optical confinement enhancement, whereas achieving comparable
absorption in Si would require centimetre-scale paths.

Once the electron-hole pair is generated within the depletion region of
width $W_i$, the built-in electric field $\mathcal{E}_0$ augmented by
the applied reverse bias $V_R$ sweeps the carriers in opposing directions:
electrons toward the n-contact and holes toward the p-contact. It is this
drift current, rather than diffusion, that constitutes the primary
photocurrent in a well-designed photodetector operated under sufficient
reverse bias.

\subsubsection{Drift Velocity and the Transit-Time Limit}
\label{subsubsec:transit_time}

The drift velocity $v_d$ of a carrier species $s \in \{e, h\}$ is,
in the low-field linear regime, proportional to the local electric field
$\mathcal{E}$ through the carrier mobility $\mu_s$:
\begin{equation}
    v_{d,s} = \mu_s \mathcal{E}.
    \label{eq:drift_linear}
\end{equation}
However, at the electric field magnitudes present in a reverse-biased
Ge p-i-n at even modest bias voltages (\SI{1}{\volt} across a
\SI{1}{\micro\metre} intrinsic region implies
$\mathcal{E} \approx \SI{e6}{\volt\per\metre}$), the velocity-field
relationship saturates. The empirical single-parameter saturation model
due to Caughey and Thomas captures this behaviour:
\begin{equation}
    v_{d,s}(\mathcal{E}) =
    \frac{\mu_{s,0}\,\mathcal{E}}
         {\left[1 + \left(\dfrac{\mu_{s,0}\,\mathcal{E}}{v_{\mathrm{sat},s}}\right)^{\beta_s}\right]^{1/\beta_s}},
    \label{eq:velocity_saturation}
\end{equation}
where $\mu_{s,0}$ is the low-field mobility, $v_{\mathrm{sat},s}$ is the
saturation velocity, and $\beta_s$ is a fitting parameter of order unity.
For germanium, the saturation velocities of electrons and holes are
$v_{\mathrm{sat},e} \approx \SI{6e4}{\metre\per\second}$ and
$v_{\mathrm{sat},h} \approx \SI{4.6e4}{\metre\per\second}$,
respectively~\cite{Yan2024}. The asymmetry between these values is
physically significant: holes transit the depletion region more slowly
than electrons, and it is therefore holes that impose the more stringent
transit-time limitation.

The transit time for a carrier species $s$ traversing a fully depleted
intrinsic layer of width $W_i$ at saturation velocity is:
\begin{equation}
    \tau_{\mathrm{tt},s} = \frac{W_i}{v_{\mathrm{sat},s}}.
    \label{eq:transit_time_species}
\end{equation}
The corresponding small-signal transit-time bandwidth limit follows from
the frequency response of the photocurrent under the assumption of uniform
carrier generation and saturated drift~\cite{SzeNg2006}:
\begin{equation}
    f_{\mathrm{tt}} = \frac{0.45\,v_{\mathrm{sat},h}}{W_i},
    \label{eq:lit_bw_tt}
\end{equation}
where the numerical prefactor 0.45 arises from the \SI{3}{\decibel} roll-off
condition for a current pulse of duration $\tau_{\mathrm{tt},h}$ and is
hole-velocity limited because $v_{\mathrm{sat},h} < v_{\mathrm{sat},e}$.
Equation~\eqref{eq:bw_tt} reveals a fundamental geometric constraint:
to achieve $f_{\mathrm{tt}} \geq \SI{100}{\giga\hertz}$, the intrinsic
layer width must satisfy
\begin{equation}
    W_i \leq \frac{0.45 \times \SI{4.6e4}{\metre\per\second}}
                  {\SI{100e9}{\per\second}} \approx \SI{207}{\nano\metre},
    \label{eq:wi_limit}
\end{equation}
a value achievable in modern CMOS-foundry Ge epitaxy but one that
simultaneously drives the junction capacitance upward as $W_i$ decreases,
creating the fundamental transit-time versus RC-delay trade-off analysed
in \cref{subsubsec:rc_bandwidth}.

\subsubsection{Junction Capacitance and the RC Bandwidth Limit}
\label{subsubsec:rc_bandwidth}

The junction capacitance $C_j$ of a uniformly depleted p-i-n structure
with cross-sectional area $A_j$ and intrinsic layer permittivity
$\varepsilon_{\mathrm{Ge}}$ is:
\begin{equation}
    C_j = \frac{\varepsilon_{\mathrm{Ge}}\,A_j}{W_i},
    \label{eq:lit_cj}
\end{equation}
where $\varepsilon_{\mathrm{Ge}} = 16.2\,\varepsilon_0$ for germanium. For a
vertical p-i-n photodetector with Ge mesa footprint $W_m \times L$ and
intrinsic layer width $W_i$, the capacitance scales linearly with the
optical mode cross-section. An increase in $L$ to improve absorption
through \cref{eq:abs_fraction} therefore increases $C_j$ proportionally.

The RC-limited bandwidth of the photodetector is determined by the time
constant $\tau_{\mathrm{RC}} = R_s C_j$, where $R_s$ is the total series
resistance of the junction, encompassing the metal-semiconductor contact
resistance, the bulk semiconductor resistance in the p- and n-type cladding
layers, and the resistance of the lateral current path from the mesa edge
to the nearest metal electrode. The \SI{3}{\decibel} RC bandwidth is:
\begin{equation}
    f_{\mathrm{RC}} = \frac{1}{2\pi\,R_s\,C_j} = \frac{W_i}{2\pi\,R_s\,\varepsilon_{\mathrm{Ge}}\,A_j}.
    \label{eq:lit_bw_rc}
\end{equation}
The combined bandwidth, accounting for both transit-time and RC
limitations under the standard root-sum-of-squares approximation valid
when the two poles are well-separated, is:
\begin{equation}
    \frac{1}{f_{3\,\mathrm{dB}}^2} \approx \frac{1}{f_{\mathrm{tt}}^2} + \frac{1}{f_{\mathrm{RC}}^2}.
    \label{eq:lit_bw_total}
\end{equation}

The interplay between \cref{eq:bw_tt,eq:bw_rc,eq:cj} defines the
central design space of high-speed photodetectors. Reducing $W_i$
improves $f_{\mathrm{tt}}$ but simultaneously increases $C_j \propto 1/W_i$
and thus degrades $f_{\mathrm{RC}}$ unless $R_s$ is proportionally
reduced or $A_j$ is decreased. Reducing $A_j$ (by narrowing the
Ge mesa or shortening the device) reduces $C_j$ and improves
$f_{\mathrm{RC}}$, but also reduces the single-pass absorbed photon
fraction $\mathcal{A}$ and hence the responsivity $\mathcal{R}$. These
opposing trends are not merely mathematical artefacts; they reflect genuine
physical constraints imposed by electromagnetism and semiconductor
transport physics. A successful high-speed photodetector design must
navigate this trade-off explicitly, identifying which limit is operative
and targeting the corresponding lever for improvement.

\subsubsection{Dark Current: Physical Origins and Decomposition}
\label{subsubsec:dark_current_physics}

In the absence of illumination, a reverse-biased p-i-n junction sustains
a residual \emph{dark current} $I_d$ that constitutes a fundamental noise
floor for the photodetection process. The dark current has two physically
distinct origins that behave differently with respect to bias, temperature,
and material quality.

The first contribution is the \emph{diffusion dark current} $I_{\mathrm{diff}}$,
arising from the thermal generation of minority carriers in the quasi-neutral
p- and n-type regions adjacent to the depletion region. These carriers
diffuse to the junction boundary, where they are swept across by the built-in
field, producing a reverse saturation current. For a one-sided abrupt junction
in which the n-type region dominates, this contribution takes the
Shockley form:
\begin{equation}
    I_{\mathrm{diff}} = A_j\,q\,\frac{D_p\,n_i^2}{L_p\,N_D},
    \label{eq:lit_jdark_diff}
\end{equation}
where $D_p$ is the minority hole diffusivity, $n_i$ is the intrinsic
carrier concentration, $L_p$ is the minority carrier diffusion length,
and $N_D$ is the donor concentration in the n-region. The $n_i^2$
dependence is critical: since $n_i \propto \exp(-E_g/2k_BT)$, the
diffusion dark current of germanium ($E_g \approx \SI{0.66}{\electronvolt}$)
is several orders of magnitude larger than that of silicon
($E_g \approx \SI{1.12}{\electronvolt}$) at room temperature. This is
the fundamental thermodynamic cost of using germanium as an absorber
material, and it motivates careful doping profile design to minimise the
volume of lightly doped germanium exposed to minority-carrier injection.

The second contribution is the \emph{generation-recombination (G-R) current}
$I_{\mathrm{GR}}$, arising from carrier generation through mid-gap trap
states (Shockley-Read-Hall centres) within the depletion region:
\begin{equation}
    I_{\mathrm{GR}} = A_j\,q\,\frac{n_i\,W_i}{2\,\tau_0},
    \label{eq:jdark_gr}
\end{equation}
where $\tau_0$ is the effective carrier lifetime determined by the trap
density and capture cross-section. The G-R current scales linearly with
$n_i$ (rather than $n_i^2$), exhibits a weaker temperature dependence than
the diffusion term, and grows with depletion width $W_i$. In high-quality
epitaxial Ge-on-Si, the G-R current frequently dominates because the
heteroepitaxial interface introduces a density of threading dislocations
with $\rho_{\mathrm{TD}} \approx 10^7$--$10^8\,\mathrm{cm}^{-2}$, each acting
as a Shockley-Read-Hall centre that reduces the effective lifetime $\tau_0$.

The total dark current is therefore:
\begin{equation}
    I_d = I_{\mathrm{diff}} + I_{\mathrm{GR}},
    \label{eq:lit_idark_decomposed}
\end{equation}
and reducing $I_d$ in a Ge-on-Si detector requires addressing both the
thermodynamic (bandgap) and materials-quality (threading dislocation density)
contributions simultaneously. The shot noise power spectral density associated
with $I_d$ is $S_{i,d} = 2qI_d$, which directly raises the noise floor and
degrades the minimum detectable optical power. For the receiver sensitivity
context of \cref{subsec:pd_gaps}, a dark current $I_d \lesssim \SI{10}{\nano\ampere}$
is a practical upper bound to avoid significant noise penalty in a
\SI{106.25}{\giga\bit\per\second} PAM-4 receiver.

\subsubsection{Responsivity and Quantum Efficiency}
\label{subsubsec:responsivity_qe}

The external quantum efficiency $\eta_{\mathrm{ext}}$ of a waveguide-integrated
photodetector quantifies the fraction of incident photons that contribute
to the measured photocurrent, accounting for all loss mechanisms:
\begin{equation}
    \eta_{\mathrm{ext}} = \eta_c\,(1-\Gamma)\,\mathcal{A}\,\eta_i,
    \label{eq:lit_qe_full}
\end{equation}
where $\eta_c$ is the optical coupling efficiency from the waveguide mode
to the Ge absorber mode, $\Gamma$ is the power reflectance at the waveguide-to-Ge
interface, $\mathcal{A} = 1 - e^{-\alpha_{\mathrm{Ge}}\Gamma_{\mathrm{opt}}L}$
is the single-pass absorption fraction for a device of length $L$ with optical
confinement factor $\Gamma_{\mathrm{opt}}$ in the Ge layer, and $\eta_i$ is the
internal quantum efficiency (the probability that an absorbed photon generates
a carrier pair that is successfully collected before recombination). In a
well-designed vertical p-i-n structure with sufficient reverse bias to deplete
the entire intrinsic Ge layer, $\eta_i \approx 1$ because the carrier sweep-out
time $\tau_{\mathrm{tt}} \ll \tau_0$.

The responsivity, defined as the ratio of photocurrent to incident optical power,
is related to $\eta_{\mathrm{ext}}$ by:
\begin{equation}
    \mathcal{R} = \frac{\eta_{\mathrm{ext}}\,q}{h\nu}
               = \frac{\eta_{\mathrm{ext}}\,\lambda}{hc},
    \label{eq:lit_responsivity}
\end{equation}
where $\lambda = \SI{1310}{\nano\metre}$, $h$ is Planck's constant, and
$c$ is the speed of light. The theoretical maximum responsivity at
\SI{1310}{\nano\metre} with $\eta_{\mathrm{ext}} = 1$ is
$\mathcal{R}_{\mathrm{max}} = \lambda q / hc \approx \SI{1.057}{\ampere\per\watt}$.
Practical vertical Ge-on-Si devices with butt-coupling achieve
$\mathcal{R} \approx 0.9$--$\SI{1.0}{\ampere\per\watt}$, corresponding
to $\eta_{\mathrm{ext}} \approx 85$--$95\%$, with the residual loss
attributed to mode mismatch at the Si-to-Ge transition and residual
reflectance at the input facet.

The relationship between $\mathcal{R}$ and device length $L$ through
$\mathcal{A}(L)$ is nonlinear and approaches saturation asymptotically.
For $\alpha_{\mathrm{Ge}}\Gamma_{\mathrm{opt}} \approx
\SI{1.5e4}{\per\metre}$ in a standard butt-coupled vertical geometry,
a device length $L = \SI{10}{\micro\metre}$ achieves
$\mathcal{A} \approx 0.78$, while doubling to $L = \SI{20}{\micro\metre}$
yields $\mathcal{A} \approx 0.95$, a diminishing return. Simultaneously,
$C_j$ scales linearly with $L$ through \cref{eq:cj}, so the responsivity
gain from lengthening the device is bought at the cost of a proportionally
larger RC time constant. This quantitative observation motivates the
search for electrode topologies that reduce $R_s$ independently of $L$,
as explored in the U-shaped anode geometry of \cref{subsec:pd_u_electrode_review}.

\subsubsection{Inductive Gain Peaking: Physical Mechanism and Design Implications}
\label{subsubsec:inductive_peaking}

The standard RC roll-off of \cref{eq:bw_rc} can be partially compensated
by introducing a series inductance $L_p$ between the photocurrent source
and the load resistance $R_L$. The resulting second-order transfer function
exhibits a controlled resonance that extends the \SI{3}{\decibel} bandwidth
beyond the unpeaked value. For a circuit consisting of a current source
$I_{\mathrm{ph}}$ in parallel with $C_j$, a series resistance $R_s$, a
series inductance $L_p$, and a load $R_L$, the voltage transfer function
has a denominator:
\begin{equation}
    H(s) = \frac{1}
    {L_p C_j\,s^2 + (R_s + R_L)C_j\,s + 1 + (R_s+R_L)/Z_0},
    \label{eq:inductive_peaking_tf}
\end{equation}
where $Z_0$ represents any shunt impedance to ground. The resonance
condition occurs when the imaginary parts of the numerator and denominator
cancel, approximately at:
\begin{equation}
    \omega_{\mathrm{res}} = \frac{1}{\sqrt{L_p C_j}},
    \label{eq:resonance_condition}
\end{equation}
and the degree of peaking is controlled by the Q-factor
$Q = \omega_{\mathrm{res}} / [\omega_{\mathrm{res}}(R_s + R_L)C_j]$.
A Butterworth (maximally flat) response is achieved at $Q = 1/\sqrt{2}$,
which sets the optimal ratio $L_p / [(R_s + R_L)^2 C_j] = 1/2$. Excessive
Q leads to an undesirable passband ripple that degrades the eye diagram
and increases inter-symbol interference in PAM-4 signalling.

The interaction between $R_s$ and the inductive peaking condition has a
direct implication for electrode engineering. Because the optimal $L_p$
scales as $(R_s + R_L)^2 C_j$, a reduction in $R_s$ (achieved, for example,
by the U-shaped anode of \cref{subsec:pd_u_electrode_review}) shifts the
optimal inductance to a smaller value that is more easily realised as a
bondwire or on-chip spiral without requiring impractically large physical
dimensions. The U-shaped topology therefore not only reduces
$\tau_{\mathrm{RC}}$ directly, but also brings the inductive peaking
condition into a more practical operating regime, providing a second-order
benefit that is not captured by the simple $f_{\mathrm{RC}}$ metric alone.

\subsubsection{Synthesis: The Coupled Design Space}
\label{subsubsec:design_space_synthesis}

The foregoing analysis identifies five coupled variables that determine
the performance of a vertical Ge-on-Si PIN photodetector:
the intrinsic layer width $W_i$, the device length $L$, the mesa
cross-sectional area $A_j = W_m \times L$, the series resistance $R_s$,
and the optional series inductance $L_p$. These variables are constrained
by four performance objectives: $f_{3\,\mathrm{dB}} \geq \SI{40}{\giga\hertz}$
for 400GBASE-DR4 compatibility, $\mathcal{R} \geq \SI{0.85}{\ampere\per\watt}$
for adequate receiver sensitivity, $I_d \leq \SI{10}{\nano\ampere}$ for
acceptable noise, and CMOS-process compatibility which constrains the
minimum $W_i$ and the available epitaxial germanium quality.

No single variable can be independently optimised without perturbing
at least one other objective. Reducing $W_i$ improves $f_{\mathrm{tt}}$
but worsens $C_j$ and may increase G-R dark current through
\cref{eq:jdark_gr} if $W_i$ enters the regime of strong field
non-uniformity near the heterojunction interface. Reducing $A_j$
improves $f_{\mathrm{RC}}$ but limits $\mathcal{R}$ through $\mathcal{A}$.
Reducing $R_s$ through electrode engineering improves $f_{\mathrm{RC}}$
without a direct penalty to either $\mathcal{R}$ or $I_d$, making it the
structurally cleanest lever available within a fixed process design kit.
This observation provides the physical and quantitative foundation for
the electrode topology choices examined in subsequent sections of this chapter.
